\begin{definition}[The Euclidean setting on $\mathbb R^{N}$, named]
For each $N\in\mathbb N$, let $\mathbb G_E(N)$ denote one sub-elliptic setting
(\noderef{CarnotSetting}) witnessing \noderef{EuclideanCarnotSettingExists} at $N$, chosen
once and for all. Writing $n_E,m_E,r_E,*_E,\|\cdot\|_E,\pi_1^E,\exp_E,\sigma_E$ for the data
of $\mathbb G_E(N)$, the choice therefore satisfies the seven identities that
\noderef{EuclideanCarnotSettingExists} supplies:
\[
 n_E=N,\quad m_E=N,\quad r_E=1,\quad x*_Ey=x+y,\quad \|x\|_E=|x|,\quad
 \exp_E\big(\pi_1^E(x)\big)=x,\quad \sigma_E(x)_i=\exp_E(e_i)
\]
for all $x,y\in\mathbb R^{n_E}$ and every $i$, where $e_i$ is the $i$-th standard basis
vector of $\mathbb R^{N}$.

Those seven identities, together with the fields that every sub-elliptic setting carries
(\noderef{CarnotSetting}), are all that this definition makes available about
$\mathbb G_E(N)$: it fixes a witness and asserts nothing further about it. In particular
$\pi_1^E$ is \emph{not} the identity map here, and no consumer may treat it as one; the
displayed identities relate $\pi_1^E$ and $\exp_E$ only as mutually inverse maps between
$\mathbb R^{n_E}$ and $\mathbb R^{N}$.

The purpose of fixing a witness, rather than invoking
\noderef{EuclideanCarnotSettingExists} where it is needed, is that a setting obtained inside
a proof cannot be named in a statement. Case (iii) of Theorem~1.3 descends to $\mathbb R^{m}$
and applies the cited results of the literature \emph{there}, at the Euclidean setting; those
results are asserted only at the setting carried as a parameter of the cited-results package,
so a hypothesis asserting them at the Euclidean setting has to be able to refer to that
setting by name. With $\mathbb G_E(m)$ named, the hypothesis can be written; without it, it cannot.
\end{definition}
